\section{Discussion}
\label{sec:discussion}

\subsection{Why Do Metadata-Mimicking Attacks Always Succeed?}

The most striking finding is that three subtle attacks maintain $\sim$100\% \asr across all experimental conditions.
These attacks share a common strategy: they frame adversarial instructions as editorial notes, formatting guidelines, or author preferences---document metadata that the model treats as legitimate processing instructions rather than content to summarize.

We hypothesize that this vulnerability arises from a fundamental ambiguity in the model's task: when asked to ``summarize a document,'' the model must distinguish between \emph{content} (to be summarized) and \emph{meta-instructions} (to be followed).
Attacks that exploit this ambiguity by disguising instructions as meta-instructions bypass the model's safety training, which focuses on detecting overt override attempts.

This interpretation is consistent with the finding that \gptfull resists direct overrides (0\% \asr) but remains fully vulnerable to metadata-mimicking attacks (100\% \asr): safety training addresses the \emph{form} of the attack (explicit instruction override) rather than the \emph{mechanism} (instruction-data confusion).

\subsection{The 2{,}000-Token Plateau}

For length-sensitive attacks, \asr drops sharply between 500 and 2{,}000 tokens but then plateaus through 32{,}000 tokens.
This plateau suggests that the model's effective attention window for deciding whether to follow an embedded instruction is limited to roughly 2{,}000 tokens of surrounding context.
Beyond this window, additional filler text provides no further dilution benefit.

This finding has practical implications: if length-based defenses are employed (\eg padding documents with benign text), they saturate quickly and provide diminishing returns beyond a modest threshold.

\subsection{Recency Bias and the End Position}

Direct attacks placed at the end of the document achieve 61.4\% \asr versus 21.7\% at all other positions.
This asymmetric position effect---concentrated at the end rather than following the symmetric U-shape predicted by ``Lost in the Middle''~\citep{liu2023lost}---suggests that for instruction-following (as opposed to factual retrieval), the model exhibits a strong recency bias: the last content it processes exerts disproportionate influence on its behavior.

This has important implications for document processing pipelines: content at the end of a document is the most dangerous position for adversarial injection, and defenses should prioritize sanitizing document tails.

\subsection{Limitations}
\label{sec:limitations}

\para{Evaluation method.}
Our pattern-matching evaluation is conservative and may miss cases of partial compliance where the model follows the spirit but not the exact letter of the injection.
An LLM-as-judge evaluation~\citep{yi2024llmjudge} would capture a richer spectrum of compliance behaviors.

\para{Attack diversity.}
Our 10 attacks, while covering direct and subtle categories, do not exhaust the space of possible injection techniques.
Notably, we do not test encoded attacks (base64, Unicode), multi-turn attacks, or gradient-optimized adversarial suffixes~\citep{pasquini2024neural}.

\para{Task scope.}
We test only document summarization.
Other tasks (\eg question answering, code generation, data extraction) may show different vulnerability patterns, particularly tasks where the model must extract and follow specific instructions from the document.

\para{Model coverage.}
We test two models from a single provider (OpenAI).
While preliminary testing showed that Claude models resist all 5 direct attacks (0\% \asr), we did not conduct a full evaluation with subtle attacks.
The three vulnerability classes we identify may not generalize across all model families.

\para{Token counting.}
Our document generator uses a character-based approximation (4 characters per token) rather than exact tokenizer-based counting.
While this introduces noise in the precise length of each document, it does not affect the relative ordering of length conditions.

\para{Deterministic setting.}
All experiments use temperature 0.0.
Real deployments often use $\text{temperature} > 0$, which introduces stochastic variation that could affect both attack success and defense reliability.
